\documentclass[compress, xcolor={svgnames}, aspectratio=169, smaller]{beamer}
\newcommand{\DirectoryPath}{../../Common} 
\usepackage{\DirectoryPath/slidespreamble}
\setbeameroption{show notes}
% \setbeameroption{hide notes}

\title{Introduction to Deep Learning}
\subtitle{Session 1: Introduction}
\author{Magda Gregorová - \href{mailto:magda.gregorova@thws.de}{magda.gregorova@thws.de}}
\institute{
\includegraphics[width=0.3\textwidth]{thws-logo_vert_en_orange-rgb.png}}
\hypersetup{
pdftitle={Introduction to Deep Learning},
pdfauthor={Magda Gregorová},
pdfkeywords={organization, plan, exam}
}

\date{{\color{skyblue1}\footnotesize Licensed according to 
\href{https://creativecommons.org/licenses/by-sa/4.0}{\color{skyblue1}CC-BY-SA 4.0}}}

\begin{document}

\begin{frame}[plain]
    \titlepage 
\end{frame}

\begin{frame}{What is Deep Learning?}

\begin{center}
{\large \textit{``Pieces of code that do stuff and nobody fully understands how.''}}
\end{center}

\vskip 2em

\structure{What do \textit{you} think Deep Learning is?}
\begin{itemize}
    \item Take a moment to write down your answer
    \item We will come back to this at the end of the course
\end{itemize}

\end{frame}

\begin{frame}{Stuff - Image Classification}

\structure{Task: Image Classification}
\begin{itemize}
    \item Given an image, assign it one label from a fixed set of categories
    \item Input: raw pixels \quad $\rightarrow$ \quad Output: a class label
\end{itemize}

\vskip 1em

\structure{CIFAR-10: A Benchmark Dataset}
\begin{itemize}
    \item 60,000 colour images of $32 \times 32$ pixels
    \item 10 classes, 6,000 images per class
\end{itemize}

\vskip 1em


\begin{center}
\foreach \cls in {airplane, automobile, bird, cat, deer, dog, frog, horse, ship, truck}{%
    \begin{minipage}{0.085\textwidth}
        \centering
        \includegraphics[width=\textwidth]{cifar10/\cls.png}\\
        {\tiny \texttt{\cls}}
    \end{minipage}\hspace{2pt}
}
\end{center}

\end{frame}

\begin{frame}{The Code That Does The Stuff}

\structure{Challenge}
\begin{itemize}
    \item Classify CIFAR-10 images into 10 categories
    \item From raw pixels to correct labels - \alert{automatically}
\end{itemize}

\vskip 1em

\structure{Approach}
\begin{itemize}
    \item Ask an AI to write the code for us
    \item \textbf{Live demo:} prompting Claude to generate a PyTorch classifier
\end{itemize}

\vskip 2em

\begin{center}
{\large \textit{Let's see if we can get a machine to learn\ldots}}
\end{center}

\end{frame}

\begin{frame}{Deep Learning is Easy}

% \begin{center}
% {\large \textit{We prompted LLM. We ran a script.\\A machine learned to see.}}
% \end{center}

\structure{Anybody can do this today}
\begin{itemize}
    \item No PhD required
    \item No coding or math needed
    \item Just prompt an LLM and run the code
\end{itemize}

\vskip 3em

\begin{center}
\colorbox{structure.fg!15}{%
    \parbox{0.85\textwidth}{%
        \centering\vskip 0.5em
        {\large \textbf{So why are you here?}}\\[0.5em]
        {\normalsize Because knowing it works is not the same as knowing \alert{WHY} it works.}\\[0.3em]
        {\normalsize Only through \alert{understanding} can you build something new, fix what breaks, and push beyond what already exists.}\\[0.3em]
        % {\normalsize The goal is not to prompt better --- it is to think deeper.}
        % {\normalsize This course is about opening the black box.}
        \vskip 0.5em
    }
}
\end{center}

\end{frame}


\begin{frame}{Lecture Overview}

\tableofcontents

\end{frame}


\section{What is Deep Learning?}

\begin{frame}{What is Deep Learning?}

\vskip 1em

\structure{Family of Techniques}
\begin{itemize}
    \item \textbf{Artificial Intelligence} - machines that exhibit intelligent behaviour
    \item \textbf{Machine Learning} - machines that learn from data rather than explicit rules
    \item \textbf{Deep Learning} - machine learning using deep neural networks
    \item \textbf{GenAI} - LLMs \& image/ video/ music/ and other generation
\end{itemize}

\vskip 0.5em

\begin{center}
\begin{tikzpicture}
    % AI circle
    \draw[fill=thwspetrol!15, draw=thwspetrol, thick] (0,0) circle (2cm);
    \node[thwspetrol, font=\small\bfseries] at (0, 1.8) {AI};

    % ML circle
    \draw[fill=thwspetrol!35, draw=thwspetrol!80, thick] (0,-0.3) circle (1.6cm);
    \node[thwspetrol!90, font=\small\bfseries] at (0, 1.1) {ML};

    % DL circle
    \draw[fill=thwspetrol!60, draw=thwspetrol!90, thick] (0,-0.6) circle (1.2cm);
    \node[white, font=\small\bfseries] at (0,0.4) {DL};

    % LLMs circle
    \draw[fill=thwspetrol!90, draw=thwspetrol, thick] (0,-1) circle (0.7cm);
    \node[white, font=\small\bfseries] at (0,-1) {GenAI};

    % Callout arrow and label
    \draw[->, thick, orange] (1.9, 1.) -- (0.55, -0.95);
    \node[anchor=west, orange, font=\small\bfseries, align=left] at (2., 1) {``AI'' today};

\end{tikzpicture}
\end{center}

\end{frame}

\begin{frame}{Why Now? - Four Driving Forces}

\vskip 2em

\begin{columns}
    \begin{column}{0.45\textwidth}
        \structure{Data}
        \begin{itemize}
            \item Internet accidental datasource
            \item Benchmarks: ImageNet, GLUE
            \item Unlabelled data is data too
        \end{itemize}

        \vskip 1em

        \structure{Algorithms}
        \begin{itemize}
            \item MLP, CNN, RNN, Backprop
            \item ReLU, Batch norm, Dropout
            \item ResNets, Attention, Transformers
        \end{itemize}
    \end{column}
    \begin{column}{0.5\textwidth}
        \structure{Compute}
        \begin{itemize}
            \item GPUs, CUDA, Custom Chips
            \item AWS, Microsoft Azure, Google Cloud 
            \item Google Colab, Kaggle, Lightning.ai
        \end{itemize}

        \vskip 1em

        \structure{Community}
        \begin{itemize}
            \item PyTorch, TensorFlow, JAX
            \item Papers with code: Kaggle, HuggingFace
            \item Open Research, Conferences, arXiv
        \end{itemize}
    \end{column}
\end{columns}

\vskip 2em

\begin{center}
\colorbox{structure.fg!15}{%
    \parbox{0.85\textwidth}{%
        \centering\vskip 0.4em
        {\normalsize None of these alone was sufficient - together they triggered a revolution.}
        \vskip 0.4em
    }
}
\end{center}

\end{frame}

\begin{frame}{Data - Scale Revolution}

\begin{center}
\structure{Image Dataset Sizes Over Time}\\[0.5em]
\begin{tikzpicture}
    % Primary axis - number of images (log scale, left)
    \begin{semilogyaxis}[
        width=0.85\textwidth,
        height=6cm,
        xlabel={\small Year},
        ylabel={\small Number of Images (log scale)},
        xlabel style={font=\small},
        ylabel style={font=\small, color=thwspetrol},
        yticklabel style={font=\tiny, color=thwspetrol},
        xmin=1995, xmax=2027,
        ymin=1e3, ymax=1e10,
        ytick={1e3, 1e4, 1e5, 1e6, 1e7, 1e8, 1e9, 1e10},
        xtick={1998, 2003, 2009, 2014, 2017, 2021, 2023},
        xticklabel style={font=\tiny, anchor=north},
        grid=major,
        grid style={dashed, gray!30},
        axis y line*=left,
        axis x line*=bottom,
        name=mainaxis,
    ]

    \addplot[thwspetrol, thick, mark=*, mark size=2.5pt] coordinates {
        (1998,  70000)
        (2003,  9000)
        (2009,  1200000)
        (2014,  330000)
        (2017,  14000000)
        (2021,  400000000)
        (2023,  5000000000)
    };

    % MNIST
    \draw[gray, thin] (axis cs:1998, 70000) -- (axis cs:1998, 500000);
    \node[font=\scriptsize, anchor=south] at (axis cs:1998, 500000) {MNIST (70K)};

    % Caltech
    \draw[gray, thin] (axis cs:2003, 9000) -- (axis cs:2002, 50000);
    \node[font=\scriptsize, anchor=south] at (axis cs:2002, 50000) {Caltech-101 (9K)};

    % ImageNet
    \draw[gray, thin] (axis cs:2009, 1200000) -- (axis cs:2008, 8000000);
    \node[font=\scriptsize, anchor=south] at (axis cs:2008, 8000000) {ImageNet (1.2M)};

    % COCO
    \draw[gray, thin] (axis cs:2014, 330000) -- (axis cs:2013, 80000);
    \node[font=\scriptsize, anchor=north] at (axis cs:2013, 80000) {COCO (330K)};

    % JFT
    \draw[gray, thin] (axis cs:2017, 14000000) -- (axis cs:2018, 3000000);
    \node[font=\scriptsize, anchor=north] at (axis cs:2018, 3000000) {JFT (14M)};

    % LAION-400M
    \draw[gray, thin] (axis cs:2021, 400000000) -- (axis cs:2020, 1500000000);
    \node[font=\scriptsize, anchor=south] at (axis cs:2020, 1500000000) {LAION-400M};

    % LAION-5B
    \draw[gray, thin] (axis cs:2023, 5000000000) -- (axis cs:2024, 1500000000);
    \node[font=\scriptsize, anchor=north] at (axis cs:2024, 1500000000) {LAION-5B (5B)};

    \end{semilogyaxis}

    % Secondary axis - image side size (linear scale, right)
    % MNIST:       28px
    % Caltech-101: 200px
    % ImageNet:    224px
    % COCO:        480px
    % JFT:         224px
    % LAION-400M:  512px
    % LAION-5B:    1024px
    \begin{axis}[
        width=0.85\textwidth,
        height=6cm,
        xmin=1995, xmax=2027,
        ymin=0, ymax=1400,
        ylabel={\small Image Side Size (px)},
        ylabel style={font=\small, color=thwsorange},
        yticklabel style={font=\tiny, color=thwsorange},
        ytick={0, 200, 400, 600, 800, 1000, 1200, 1400},
        yticklabels={0, 200, 400, 600, 800, 1000, 1200, 1400},
        axis y line*=right,
        axis x line=none,
        xtick=\empty,
    ]

    \addplot[thwsorange, thick, dashed, mark=square*, mark size=2pt] coordinates {
        (1998,  28)
        (2003,  200)
        (2009,  224)
        (2014,  480)
        (2017,  224)
        (2021,  512)
        (2023,  1024)
    };

    \end{axis}

\end{tikzpicture}
\end{center}

\end{frame}


\begin{frame}{Compute \& Progress}

\begin{columns}
    \begin{column}{0.48\textwidth}
        \begin{tikzpicture}
            \begin{axis}[
                width=\textwidth,
                height=5cm,
                xlabel={\small Year},
                ylabel={\small TFLOPS (FP32)},
                xlabel style={font=\small},
                ylabel style={font=\small, color=thwspetrol},
                yticklabel style={font=\tiny, color=thwspetrol},
                xticklabel style={font=\tiny, anchor=north},
                xmin=2010, xmax=2026,
                ymin=0, ymax=2400,
                xtick={2012, 2015, 2018, 2020, 2022, 2024},
                ytick={0, 400, 800, 1200, 1600, 2000, 2400},
                grid=major,
                grid style={dashed, gray!30},
title={\structure{GPU Performance (TFLOPS)}},
        axis y line*=left,
        axis x line*=bottom,
            ]
            % GTX 680 (2012):   3.1
            % GTX 980 (2014):   4.6
            % GTX 1080 (2016):  10.6
            % V100 (2018):      14.1
            % A100 (2020):      77.6
            % RTX 4090 (2022):  165
            % H100 (2022):      756
            % H200 (2024):      1979
            \addplot[thwspetrol, thick, mark=*, mark size=2pt] coordinates {
                (2012,   3)
                (2014,   5)
                (2016,  11)
                (2018,  14)
                (2020,  78)
                (2022, 756)
                (2024, 1979)
            };
            % Labels
            \draw[gray, thin] (axis cs:2012, 3) -- (axis cs:2013, 200);
            \node[font=\scriptsize, anchor=south] at (axis cs:2013, 200) {GTX 680};

            \draw[gray, thin] (axis cs:2018, 14) -- (axis cs:2017, 250);
            \node[font=\scriptsize, anchor=south] at (axis cs:2017, 250) {V100};

            \draw[gray, thin] (axis cs:2020, 78) -- (axis cs:2019, 450);
            \node[font=\scriptsize, anchor=south] at (axis cs:2019, 450) {A100};

            \draw[gray, thin] (axis cs:2022, 756) -- (axis cs:2021, 1000);
            \node[font=\scriptsize, anchor=south] at (axis cs:2021, 1000) {H100};

            \draw[gray, thin] (axis cs:2024, 1979) -- (axis cs:2023, 2100);
            \node[font=\scriptsize, anchor=south] at (axis cs:2023, 2100) {H200};

            \end{axis}
        \end{tikzpicture}
    \end{column}
    \begin{column}{0.48\textwidth}
        \begin{tikzpicture}
            \begin{axis}[
                width=\textwidth,
                height=5cm,
                xlabel={\small Year},
                ylabel={\small Top-5 Accuracy (\%)},
                xlabel style={font=\small},
                ylabel style={font=\small, color=thwsorange},
                yticklabel style={font=\tiny, color=thwsorange},
                xticklabel style={font=\tiny, anchor=north},
                xmin=2010, xmax=2026,
                ymin=70, ymax=105,
                xtick={2012, 2015, 2018, 2020, 2022, 2024},
                ytick={70, 70, 80, 90, 100},
                grid=major,
                grid style={dashed, gray!30},
title={\structure{ImageNet Top-5 Accuracy (\%)}},
        axis y line*=left,
        axis x line*=bottom,
            ]
            % Pre-DL (2011):    74.3
            % AlexNet (2012):   84.7
            % VGG (2014):       92.7
            % ResNet (2015):    96.4
            % SENet (2017):     97.7
            % EfficientNet (2019): 98.7
            % ViT (2021):       99.0
            % CoAtNet (2022):   99.2
            \addplot[thwsorange, thick, mark=*, mark size=2pt] coordinates {
                (2011, 74.3)
                (2012, 84.7)
                (2014, 92.7)
                (2015, 96.4)
                (2017, 97.7)
                (2019, 98.7)
                (2021, 99.0)
                (2022, 99.2)
            };

            % Human performance line
            \addplot[red, dashed, thick] coordinates {
                (2010, 95.0)
                (2026, 95.0)
            };
            \node[font=\scriptsize, color=red, anchor=west] at (axis cs:2020, 96.2) {Human ($\approx$95\%)};

            % Labels
            \draw[gray, thin] (axis cs:2011, 74.3) -- (axis cs:2012, 72);
            \node[font=\scriptsize, anchor=west] at (axis cs:2012, 72) {SVM};

            \draw[gray, thin] (axis cs:2012, 84.7) -- (axis cs:2013, 80);
            \node[font=\scriptsize, anchor=north west] at (axis cs:2013, 80) {AlexNet};

            \draw[gray, thin] (axis cs:2014, 92.7) -- (axis cs:2015, 90);
            \node[font=\scriptsize, anchor=north] at (axis cs:2015, 90) {VGG};

            \draw[gray, thin] (axis cs:2015, 96.4) -- (axis cs:2014, 100);
            \node[font=\scriptsize, anchor=south] at (axis cs:2014, 100) {ResNet};

            \draw[gray, thin] (axis cs:2019, 98.7) -- (axis cs:2018, 101);
            \node[font=\scriptsize, anchor=south] at (axis cs:2018, 101) {EfficientNet};

            \draw[gray, thin] (axis cs:2021, 99.0) -- (axis cs:2022, 101);
            \node[font=\scriptsize, anchor=south] at (axis cs:2022, 101) {ViT};

            \end{axis}
        \end{tikzpicture}
    \end{column}
\end{columns}

\end{frame}

\section{Brief History}

\begin{frame}{Brief History of Deep Learning}

\begin{tikzpicture}
    \begin{axis}[
        width=\textwidth,
        height=6cm,
        xlabel={\small Year},
        xlabel style={font=\small},
        xmin=1955, xmax=2028,
        ymin=0, ymax=1,
        axis y line=none,
        axis x line=bottom,
        xtick={1958, 1969, 1986, 1998, 2006, 2012, 2017, 2022},
        xticklabel style={font=\tiny, rotate=0, anchor=north},
        grid=none,
        clip=false,
    ]

    % Timeline line
    \addplot[gray, thick] coordinates {(1955, 0.5) (2028, 0.5)};

    % Shaded AI winters
    \fill[structure.fg!5] (axis cs:1974, 0) rectangle (axis cs:1980, 1);
    \fill[structure.fg!5] (axis cs:1988, 0) rectangle (axis cs:1993, 1);
    \fill[structure.fg!5] (axis cs:1998, 0) rectangle (axis cs:2005, 1);
    \node[font=\scriptsize, structure.fg, rotate=90] at (axis cs:1977, 0.5) {AI Winter};
    \node[font=\scriptsize, structure.fg, rotate=90] at (axis cs:1990.5, 0.5) {AI Winter};
    \node[font=\scriptsize, structure.fg, rotate=90] at (axis cs:2002, 0.5) {Cold Spell};

    % Events - alternating above and below the line
    \draw[thwspetrol, thick] (axis cs:1958, 0.5) -- (axis cs:1958, 0.75);
    \node[font=\scriptsize, thwspetrol, anchor=south, align=center] at (axis cs:1958, 0.75) {Perceptron};
    \draw[thwspetrol, thick] (axis cs:1969, 0.5) -- (axis cs:1969, 0.25);
    \node[font=\scriptsize, thwspetrol, anchor=north, align=center] at (axis cs:1969, 0.25) {XOR Problem};    

    \draw[thwspetrol, thick] (axis cs:1986, 0.5) -- (axis cs:1986, 0.75);
    \node[font=\scriptsize, thwspetrol, anchor=south, align=center] at (axis cs:1986, 0.75) {Backprop};
    \draw[thwspetrol, thick] (axis cs:1998, 0.5) -- (axis cs:1998, 0.25);
    \node[font=\scriptsize, thwspetrol, anchor=north, align=center] at (axis cs:1998, 0.25) {LeNet};

    \draw[thwspetrol, thick] (axis cs:2006, 0.5) -- (axis cs:2006, 0.75);
    \node[font=\scriptsize, thwspetrol, anchor=south, align=center] at (axis cs:2006, 0.75) {Deep Belief Nets};
    \draw[thwspetrol, thick] (axis cs:2012, 0.5) -- (axis cs:2012, 0.25);
    \node[font=\scriptsize, thwspetrol, anchor=north, align=center] at (axis cs:2012, 0.25) {AlexNet};  

    \draw[thwspetrol, thick] (axis cs:2017, 0.5) -- (axis cs:2017, 0.75);
    \node[font=\scriptsize, thwspetrol, anchor=south, align=center] at (axis cs:2017, 0.75) {Transformers};
    \draw[thwspetrol, thick] (axis cs:2022, 0.5) -- (axis cs:2022, 0.25);
    \node[font=\scriptsize, thwspetrol, anchor=north, align=center] at (axis cs:2022, 0.25) {ChatGPT}; 



    \end{axis}
\end{tikzpicture}

\end{frame}

\section{Applications}

\begin{frame}{What Can Deep Learning Do?}

\begin{columns}
    \begin{column}{0.48\textwidth}
        \structure{Vision}
        \begin{itemize}
            \item Image classification \& object detection
            \item Medical imaging - tumour detection, X-ray analysis
            \item Autonomous driving
            \item Face recognition
        \end{itemize}

        \vskip 1em

        \structure{Language}
        \begin{itemize}
            \item Machine translation
            \item Question answering - ChatGPT, Claude
            \item Code generation
            \item Document summarisation
        \end{itemize}
    \end{column}
    \begin{column}{0.48\textwidth}
        \structure{Generative Models}
        \begin{itemize}
            \item Image generation - Stable Diffusion, DALL-E
            \item Video generation
            \item Music and speech synthesis
            \item Drug and protein design
        \end{itemize}

        \vskip 1em

        \structure{Beyond}
        \begin{itemize}
            \item Game playing - AlphaGo, AlphaFold
            \item Robotics and control
            \item Scientific discovery
            \item Recommendation systems
        \end{itemize}
    \end{column}
\end{columns}

\end{frame}

\section{The Deep Learning Pipeline}
\section{Course Roadmap \& Logistics}

\end{document}


% \begin{frame}{About me}

%     \begin{columns}
%         % Left column for the image
%         \begin{column}{0.7\textwidth}
%             \structure{\large Magda Gregorová}
%         \end{column}
        
%         % Right column for the title and content
%         \begin{column}{0.1\textwidth}
%             \includegraphics[width=\textwidth]{MG.png}
%         \end{column}
%     \end{columns}

%     % \vspace{0.5cm} % Add vertical space

%     \begin{itemize}
%         \item 2002 MSc Statistics, University of Economics, Prague, CZ
%         \item 2001-07 Monetary statistics, Czech National Bank (ECB, IMF)
%         \item 2008-12 Forecaster, Eurocontrol, Bruxelles
%         \item 2012-13 Senior statistician, Expedia, Geneva
%         \item 2013-18 PhD Machine learning, University of Geneva, CH
%         \item 2018-21 Senior researcher, DMML Group of Geneva, CH
%         \item \structure{2021- Prof for representation learning THWS/FIW}
%         \item 2022-24 Head of CAIRO, 2023- Member of the faculty board
%     \end{itemize}

% \end{frame}